\section{Transport calculation for the four-copy witness}\label{sec:transport}

We now explain the rank jump behind \eqref{eq:q13-values}.  Fix the six copy
permutations
\[
\boldsymbol\sigma=(1,(03)(12),(23),(02)(13),(01)(23),(021)).
\]
Let \(V\) be the three-dimensional message space and let
\(g_i:V\to k\) be the six coordinate covectors.  Form a bipartite
four-sheeted cover with vertices \(L_a,R_b\), \(0\leq a,b<4\), and a
color-\(i\) edge from \(L_a\) to \(R_{\sigma_i(a)}\).  Put \(V\) at each
vertex and \(k\) at each edge, using \(g_{\eta(i)}\) as both restriction
maps for a coordinate assignment \(\eta\in S_6\).  This is a transport
system: each edge equation requires the vertex data at its two ends to agree
after applying the corresponding coordinate functional.  A solution is a tuple
\((x_a,y_b)\in V^8\) satisfying
\begin{equation}\label{eq:transport-section}
 g_{\eta(i)}(x_a-y_{\sigma_i(a)})=0 .
\end{equation}
on all \(24\) edges.  The three-dimensional constant-section space is
always present.  Gauging it away leaves the \(24\)-by-\(21\) matching
matrix in \eqref{eq:permutation-contraction}.

Let \(U=k^4/k(1,1,1,1)\), and let \(\rho_i\) be the action of
\(\sigma_i\) on \(U\).  For each \(\eta\), the MDS property makes
\(g_{\eta(3)},g_{\eta(4)},g_{\eta(5)}\) a basis of \(V^*\).  Define
\(Q_\eta\) by
\begin{equation}\label{eq:transport-systematic}
 g_{\eta(i)}=\sum_{k=0}^2(Q_\eta)_{ik}g_{\eta(3+k)}
 \qquad(0\leq i<3).
\end{equation}
Thus \(Q_\eta\) is the left block obtained by systematizing the
party-reordered generator matrix on its last three columns; no pivot can
vanish because every three coordinate covectors are independent.  Put
\(\widetilde Q_\eta=2(t-1)Q_\eta\).  This scalar is a unit on the regular
non-GRS locus.  For the identity assignment, the polynomial representative is
\[
\widetilde Q_1=\begin{pmatrix}a&b&c\\a&c&b\\d&d&d\end{pmatrix},
\quad
\begin{array}{ll}
a=-2(t-1)^2,&b=-(t^2-t+1),\\
c=-(t^2-3t+1),&d=-2(t-1).
\end{array}
\]
If \(Y_k\in U\) are the three free copy vectors, then
\eqref{eq:transport-section} reduces
to the \(9\)-by-\(9\) block transport operator
\begin{equation}\label{eq:transport-operator}
 \widetilde T_\eta
 =\bigl((\widetilde Q_\eta)_{ik}(\rho_i-\rho_{3+k})\bigr)_{i,k=0}^2.
\end{equation}
The dimension identity
\begin{equation}\label{eq:transport-rank-bridge}
 21-\rank M_{\pi\boldsymbol\sigma}(G(t))
   =9-\rank\widetilde T_\eta
\end{equation}
holds over the field \(k\), where \(\eta=\pi\) in the convention of
Section~\ref{sec:invariants}.  It follows by quotienting the
three-dimensional constant-section kernel and then eliminating the three
free systematic coordinates.  Multiplying the resulting operator by
the unit \(2(t-1)\) does not change its rank.  All steps are invertible
on the regular non-GRS six-arc locus.

The remaining calculation has an internal group symmetry.  The next lemma
reduces the \(720\) coordinate assignments to four double cosets for each
rank-drop question.  Cycles act on the coordinate labels \(1,\ldots,6\).

\begin{lemma}[Double-coset reduction]\label{lem:transport-double-cosets}
Define
\[
\begin{aligned}
H_s&=\langle(34)(56),(3546),(12)(56),(13)(24)(56)\rangle,&
K_s&=\langle(36),(23)(46),(12)(45)\rangle,\\
H_a&=\langle(56),(34),(35)(46),(12),(13)(24)\rangle,&
K_a&=\langle(36),(23)(46),(15)\rangle.
\end{aligned}
\]
Then \((|H_s|,|K_s|)=(24,48)\) and
\((|H_a|,|K_a|)=(48,16)\).  The signed double cosets
\[
 H_s\backslash S_6/K_s
 \quad\text{have representatives}\quad
 1,(56),(456),(235)
\]
have sizes \(192,192,288,48\), respectively.  The axial double cosets
\[
 H_a\backslash S_6/K_a
 \quad\text{have representatives}\quad
 1,(456),(2345),(235)
\]
have sizes \(384,192,96,48\), respectively.

At a generic point of \(3B-2A=0\), the reduced transport operator has rank
eight exactly on \(H_sK_s\).  At a generic point of \(3B+2A=0\), it has
rank eight exactly on \(H_s(56)K_s\).  At a generic point of
\(B^2-2A^2=0\), it has rank eight exactly on
\(H_a(2345)K_a\).  On all other double cosets its rank is nine.
\end{lemma}

\begin{proof}
The displayed generators act on \(\widetilde T_\eta\) by invertible row and
column operations on the indicated divisor, so rank is constant on the
corresponding double cosets.  Multiplying the generators gives subgroup
orders \(24,48,48,16\).  Their intersections at the listed representatives
have orders \(6,6,4,24\) in the signed case and \(2,4,8,16\) in the axial
case.  The formula
\[
 |HrK|=|H||K|/|H\cap rKr^{-1}|
\]
therefore gives the displayed sizes; in each case the four sizes sum to
\(|S_6|=720\), so the lists are exhaustive.

It remains to evaluate one \(9\)-by-\(9\) operator per double coset.  Using
\(a=b+c\), \(d^2=-2a\), \(A=a(c-b)\), and \(B=bc\), fraction-free
elimination gives, at the three rank-dropping representatives,
\[
\begin{aligned}
\det\widetilde T_1&=64d^3AB(3B-2A),\\
\det\widetilde T_{(56)}&=64d^3AB(3B+2A),\\
\det\widetilde T_{(2345)}&=64d^3A(B^2-2A^2).
\end{aligned}
\]
The representative determinants on the other double cosets are nonzero
modulo the corresponding factor.  A nonzero \(8\)-minor remains on each of
the three components, so every stated rank drop is exactly from nine to
eight.
\end{proof}

\begin{theorem}[Transport rank-drop locus]\label{thm:transport-divisor}
For a parameter on the regular odd non-GRS pencil, there exists a coordinate
assignment \(\eta\in S_6\) for which the transport system
\eqref{eq:transport-section} has a nonconstant
section if and only if
\[
 (B^2-2A^2)(9B^2-4A^2)=0,
\]
or equivalently
\begin{equation}\label{eq:transport-divisor}
 (z-2)(9z-4)=0.
\end{equation}
At a generic point above \(z=2\), exactly \(96\) coordinate assignments have
rank \(20\); at a generic point above \(z=4/9\), exactly \(192\) do.
Every other assignment has rank \(21\).
\end{theorem}

\begin{proof}
Lemma~\ref{lem:transport-double-cosets} shows that the only generic rank-drop
components are \(3B\pm2A=0\) and \(B^2-2A^2=0\).  The first two are the two
sheets \(B/A=\mp2/3\) above \(z=4/9\), and each supports \(192\) coordinate
assignments.  The last is \(z=2\) and supports \(96\) assignments.  The
nonzero \(8\)-minors make every reduced kernel one-dimensional.
Equation~\eqref{eq:transport-rank-bridge} turns these kernels into rank
\(20\) of the matching matrix and gives rank \(21\) everywhere else.
\end{proof}

The exceptional arithmetic belongs to this detector, not to
Theorem~\ref{thm:lu-lc-rigidity}.

\begin{corollary}[Exceptional characteristics]\label{cor:transport-characteristics}
On the regular locus, characteristic \(3\) has no rank-drop component.  In
characteristic \(5\), only the \(z=4/9\) component disappears, because it is
GRS.  In characteristic \(7\), the two components merge; exactly \(288\)
assignments then have rank \(20\), and the remaining \(432\) have rank
\(21\).  In every other odd characteristic the two components and their
generic multiplicities remain as in Theorem~\ref{thm:transport-divisor}.
\end{corollary}

\begin{proof}
The identity \(4B^2+A^2=4G^2\) specializes to
\[
\begin{aligned}
\operatorname{char}k=3:&\quad
B^2-2A^2=G^2,\qquad 9B^2-4A^2=-A^2;\\
\operatorname{char}k=5:&\quad 9B^2-4A^2=4G^2;\\
\operatorname{char}k=7:&\quad
9B^2-4A^2=2(B^2-2A^2).
\end{aligned}
\]
The regular-locus exclusions give the first two claims.  In characteristic
\(7\), the disjoint \(96\)- and \(192\)-element double cosets of
Lemma~\ref{lem:transport-double-cosets} give the stated histogram.  The only
other primes in the resultants of the representative determinants with the
excluded boundary factors are \(11,13,41\); direct reduction of the displayed
nonzero \(8\)-minors shows that these cause ramification but no further rank
drop.
\end{proof}

The paper-local certificate independently replays the representative
determinants and their reductions.

The divisor records where one scalar contraction changes rank.  It is
not an LU-orbit boundary and, by
Proposition~\ref{thm:fixed-copy-boundary}, cannot be promoted to a
generic coordinate.
